\epanalipsi{4}

\tmima{Λεξικό του κεφαλαίου}

\renewcommand{\arraystretch}{1.5}
\begin{tabular*}{\linewidth}{@{\extracolsep{\fill}} >{\raggedright\arraybackslash}p{0.28\linewidth} >{\raggedright\arraybackslash}p{0.62\linewidth}}
\toprule
\textbf{Λέξη ή τύπος} & \textbf{Με μία πρόταση}\\
\midrule
αλγεβρική παράσταση & Έκφραση με αριθμούς, γράμματα και πράξεις, π.χ. $3x+5$.\\
εξίσωση & Ισότητα με άγνωστο γράμμα, π.χ. $2x+3=11$.\\
ρίζα, λύση & Ο αριθμός που βάζω στη θέση του αγνώστου και η ισότητα ισχύει.\\
επαλήθευση & Βάζω τη λύση στην αρχική εξίσωση (και στο πρόβλημα) και ελέγχω.\\
ταυτότητα & Ισχύει για κάθε αριθμό, π.χ. $x+3=3+x$.\\
αδύνατη & Δεν ισχύει για κανέναν αριθμό, π.χ. $x+2=x+6$.\\
$x+\alpha=\beta$ και οι άλλες & $x=\beta-\alpha$, $x=\beta+\alpha$, $x=\alpha-\beta$, $x=\beta:\alpha$, $x=\beta\cdot\alpha$, $x=\alpha:\beta$.\\
βήματα επίλυσης προβλήματος & Άγνωστος · εκφράζω · εξίσωση · λύνω · επαληθεύω.\\
\bottomrule
\end{tabular*}

\tmima{Ασκήσεις επανάληψης}

\ex Γράψε με γράμματα το πενταπλάσιο ενός αριθμού ελαττωμένο κατά $3$ και τον επόμενο του διπλάσιου ενός αριθμού $\nu$.

\ex Γράψε απλούστερα την παράσταση $7x+2y-3x+5y$.

\ex Λύσε: α) $x+3{,}4=10$ \quad β) $2{,}5x=15$ \quad γ) $20-x=7{,}5$ \quad δ) $x:4=2{,}5$ \quad ε) $18:x=4$.

\ex Λύσε: α) $3x+7=34$ \quad β) $5x-2=18$ \quad γ) $40-3x=7$ \quad δ) $\frac{x}{4}+2=6$.

\ex Χαρακτήρισε ως ταυτότητα, αδύνατη ή εξίσωση με μία λύση: α) $x+5=5+x$ \quad β) $x+5=x+7$ \quad γ) $2x=0$ \quad δ) $0\cdot x=3$.

\ex \textbf{Η δεξαμενή.} Μια δεξαμενή έχει όγκο $9$ m$^3$, μήκος $2{,}5$ μ. και πλάτος $1{,}2$ μ. Βρες το ύψος της.

\ex \textbf{Το πακέτο.} Το πακέτο κινητής έχει $5$ € πάγιο και $0{,}03$ € ανά λεπτό. Ο λογαριασμός ήταν $9{,}2$ €. Πόσα λεπτά μίλησε;

\ex \textbf{Η τούρτα.} Τα $\frac{3}{4}$ των κομματιών μιας τούρτας είναι $18$. Πόσα κομμάτια έχει η τούρτα;

\ex \textbf{Οι διαδοχικοί.} Βρες τρεις διαδοχικούς φυσικούς με άθροισμα $87$.

\ex \textbf{Η μοιρασιά.} Η Αγγελική και ο Νίκος μοιράζονται $100$ €, ο Νίκος παίρνει $12$ € περισσότερα. Πόσα παίρνει ο καθένας;

\ex \textbf{Οι χειραψίες.} Σε μια συνάντηση $15$ ατόμων όλοι δίνουν το χέρι σε όλους. Πόσες χειραψίες γίνονται;

\ex \textbf{Η αμοιβή.} Ο θείος Λευτέρης δούλεψε $5$ ημέρες και πήρε προκαταβολή τα $\frac{3}{5}$ της αμοιβής του, $240$ €. Πόσο ήταν το μεροκάματο;

\ex \textbf{Ερωτήσεις αυτοαξιολόγησης (εξισώσεις).} Σωστό ή λάθος; α) Η $2x=12$ είναι εξίσωση. \quad β) Η ρίζα της $x-7=5$ είναι το $12$. \quad γ) $3x+4x=7x^2$. \quad δ) Η $x+1=x+1$ είναι ταυτότητα. \quad ε) Η $12-x=5$ δίνει $x=17$. \quad στ) Η $0\cdot x=0$ έχει μία μόνο λύση.

\ex \textbf{Ερωτήσεις αυτοαξιολόγησης (προβλήματα).} Σωστό ή λάθος; α) Το πρώτο βήμα είναι να ονομάσω τον άγνωστο. \quad β) Αρκεί να επαληθευτεί η εξίσωση, όχι το πρόβλημα. \quad γ) Κάθε πρόβλημα λύνεται με εξίσωση. \quad δ) Η $2x+3=11$ έχει λύση το $7$. \quad ε) $\nu$ άτομα κάνουν $\frac{\nu(\nu-1)}{2}$ χειραψίες.

\tmima{Ο χάρτης της επανάληψης}

Μετά τις λύσεις, δες σε ποιες ασκήσεις δυσκολεύτηκες. Ο πίνακας σου λέει ποια παράγραφο να ξαναδιαβάσεις.

\renewcommand{\arraystretch}{1.5}
\begin{tabular*}{\linewidth}{@{\extracolsep{\fill}} >{\raggedright\arraybackslash}p{0.38\linewidth} >{\raggedright\arraybackslash}p{0.52\linewidth}}
\toprule
\textbf{Αν δυσκολεύτηκες στις} & \textbf{ξαναδιάβασε τις παραγράφους}\\
\midrule
Ε4.1, Ε4.2, Ε4.3, Ε4.5, Ε4.13 & 4.1 (έννοια της εξίσωσης)\\
Ε4.4, Ε4.6, Ε4.7, Ε4.8, Ε4.14 & 4.2 (επίλυση προβλημάτων)\\
Ε4.9, Ε4.10, Ε4.11, Ε4.12 & 4.3 (παραδείγματα επίλυσης προβλημάτων)\\
\bottomrule
\end{tabular*}
